% !TEX encoding = UTF-8
% !TEX program = lualatex

\chapter{Прототип системы выявления компьютерных атак}

В данном приложении приведён исходный код прототипа системы выявления атак на основе атрибутивных метаграфов (AMGD). Реализация выполнена на языке Python 3.10 с использованием библиотек NetworkX и Neo4j.

\section{Класс AttributiveMetaGraph}

Основной класс, реализующий динамический атрибутивный метаграф:

\begin{lstlisting}[language=Python, caption={Класс AttributiveMetaGraph}]
from typing import Dict, List, Set
import networkx as nx

class AttributiveMetaGraph:
    """
    Атрибутивный метаграф для представления событий безопасности.
    Поддерживает добавление вершин, гипердуг и сопоставление с шаблонами.
    """
    def __init__(self):
        self.graph = nx.MultiDiGraph()
        self.vertex_counter = 0
        self.edge_counter = 0

    def add_vertex(self, attrs: Dict) -> str:
        """Добавляет вершину с атрибутами и возвращает её идентификатор."""
        vid = f"v_{self.vertex_counter}"
        self.graph.add_node(vid, **attrs)
        self.vertex_counter += 1
        return vid

    def add_hyperedge(self, sources: List[str], targets: List[str], attrs: Dict) -> str:
        """
        Добавляет гипердугу (событие безопасности) между множествами источников и целей.
        """
        eid = f"e_{self.edge_counter}"
        for src in sources:
            for tgt in targets:
                self.graph.add_edge(src, tgt, id=eid, **attrs)
        self.edge_counter += 1
        return eid

    def subgraph_isomorphism(self, pattern: 'AttributiveMetaGraph', theta: float = 0.75) -> bool:
        """
        Упрощённая реализация сопоставления: проверка покрытия тактик MITRE ATT&CK.
        """
        live_tactics = {
            self.graph.edges[e]['mitre_ttp']
            for e in self.graph.edges
            if 'mitre_ttp' in self.graph.edges[e]
        }
        pattern_tactics = {
            pattern.graph.edges[e]['mitre_ttp']
            for e in pattern.graph.edges
            if 'mitre_ttp' in pattern.graph.edges[e]
        }
        if not pattern_tactics:
            return False
        overlap = len(live_tactics & pattern_tactics)
        return (overlap / len(pattern_tactics)) >= theta
\end{lstlisting}

\section{Алгоритм сопоставления шаблонов}

Функция сопоставления эталонного шаблона атаки с динамическим метаграфом:

\begin{lstlisting}[language=Python, caption={Алгоритм MetaGraph-Match}]
def match_pattern(live_graph: AttributiveMetaGraph, pattern: dict) -> float:
    """
    Возвращает меру схожести от 0.0 до 1.0 между динамическим метаграфом
    и эталонным шаблоном атаки.
    
    Аргументы:
        live_graph: AttributiveMetaGraph - текущее состояние системы
        pattern: dict - шаблон атаки в формате JSON
    
    Возвращает:
        float: мера схожести (0.0 - нет совпадений, 1.0 - полное совпадение)
    """
    score = 0.0
    matched_edges = 0
    
    for p_edge in pattern['edges']:
        tactic = p_edge['attrs']['mitre_ttp']
        # Поиск соответствующих дуг в live_graph
        for edge in live_graph.graph.edges(data=True):
            if edge[2].get('mitre_ttp') == tactic:
                score += 1.0
                matched_edges += 1
                break
    
    return score / len(pattern['edges']) if pattern['edges'] else 0.0
\end{lstlisting}

\section{Пример использования}

Пример создания шаблона атаки APT29 и его сопоставления с событиями:

\begin{lstlisting}[language=Python, caption={Пример использования}]
# Шаблон атаки APT29
apt29_pattern = {
    "edges": [
        {
            "src": ["user", "file_invoice.docm"],
            "tgt": ["process_winword"],
            "attrs": {"mitre_ttp": "T1566", "tactic": "Initial Access"}
        },
        {
            "src": ["process_winword"],
            "tgt": ["process_powershell"],
            "attrs": {"mitre_ttp": "T1059.001", "tactic": "Execution"}
        },
        {
            "src": ["process_powershell"],
            "tgt": ["ip_c2"],
            "attrs": {"mitre_ttp": "T1071.001", "tactic": "Command and Control"}
        }
    ]
}

# Построение динамического метаграфа из событий телеметрии
live_graph = AttributiveMetaGraph()

# Событие 1: открытие фишингового письма
user = live_graph.add_vertex({"type": "user", "name": "Alice"})
file = live_graph.add_vertex({"type": "file", "name": "invoice.docm", "hash": "sha256:abc123"})
winword = live_graph.add_vertex({"type": "process", "name": "winword.exe"})
live_graph.add_hyperedge([user, file], [winword], {"mitre_ttp": "T1566", "timestamp": "2025-03-15T10:23:45Z"})

# Событие 2: запуск PowerShell
powershell = live_graph.add_vertex({"type": "process", "name": "powershell.exe"})
live_graph.add_hyperedge([winword], [powershell], {"mitre_ttp": "T1059.001", "timestamp": "2025-03-15T10:24:12Z"})

# Событие 3: C2-соединение
c2 = live_graph.add_vertex({"type": "ip", "name": "91.218.112.42"})
live_graph.add_hyperedge([powershell], [c2], {"mitre_ttp": "T1071.001", "timestamp": "2025-03-15T10:25:03Z"})

# Сопоставление
similarity = match_pattern(live_graph, apt29_pattern)
if similarity >= 0.75:
    print(f"Обнаружена атака APT29! Схожесть: {similarity:.2f}")
\end{lstlisting}

Данный прототип был использован в экспериментальном исследовании, описанном в главе~4, и показал высокую эффективность при выявлении многоэтапных атак.